\subsection{The four residual links for the rank-four case}
\label{sec:count-three-chain}

The tensor now supplies a point in ordinary support. To reach the target we construct residual intermediate representations over a common finite field, so successive comparisons refer to the same local characters.

\begin{proposition}[The residual path over one finite field]\label{prop:count-three-residual-path}
Use the tower of Proposition~\ref{prop:count-three-fitted-layers},
the normalized tensor of Proposition~\ref{prop:count-three-tensor-endpoint}
and the support conditions of Lemma~\ref{hyp:count-three-path-endpoints}.
In the notation for the two bridge points \(P_3,P_4\), the common intermediate datum \(P_2\), target \(P_0\) and seed \(S_F\), the four residual lines carrying the ordered flags form the path 
\[
 P_0\longrightarrow P_3\longrightarrow P_4\longrightarrow P_2\longrightarrow S_F.
\]
 All lines and the good endpoints can be
defined over a single finite residue field after the indicated
coefficient extension.
\end{proposition}

\begin{proof}
Use the ordered frame on the tame field \(F_0\) from
Section~\ref{sec:count-three-tower}. Write \(d_i\) for
the four diagonal residual characters, and put
\(A=d_1/d_0\), \(B=d_2/d_0\), so \(d_3/d_0=AB\).
The roots of these three ratios are
\(-w_\alpha,-w_\Psi,-w_\vartheta\). Each of the
ratios is nontrivial at the normalizer. A residual
datum comprises its global matrices and one lower
unipotent symplectic frame \(N_v\) at each marked
place. The conjugated local matrices are upper
triangular with prescribed diagonal throughout
the chain.

\begin{established}[The support of the normalized target and the tensor seed]\label{hyp:count-three-path-endpoints}
On the fitted tower of Proposition~\ref{prop:count-three-fitted-layers}, write for off-diagonal entries:
\[
 \begin{split}
 S_+&=\{w_\beta,-w_\alpha,-w_\Psi,-w_\vartheta\},\\
 S_-&=\{-w_\beta,-w_\alpha,-w_\Psi,-w_\vartheta\},\\
 S_0&=S_+\cap S_- .
 \end{split}
\]
Write \(\mathcal A\) for the residual data with
global support in
\(\{-w_\beta,-w_\Psi,-w_\vartheta\}\) and
\[
 (N_v)_{10}\ne0
\]
at every marked place.
Here support refers to entries that are allowed to
be nonzero. The normalized rational target restricts
to a datum \(P_0\) with global support in \(S_+\) and frame support in
\(S_0\). Its two possible initial conditions are
\begin{enumerate}
\item the global roots \(-w_\Psi\) and \(w_\beta\) occur,
\item the global support lies in
\[
 S_+\setminus\{-w_\Psi\},
\]
all the frame
entries \((N_v)_{20}\) are nonzero and
the global root \(-w_\alpha\) occurs.
\end{enumerate}


The input from the tensor construction is one more
datum \(S\) on \(F_0\), with global support in \(S_-\).
It can be ramified at one extra split rational prime
\(w\) with \(w\equiv1\pmod\ell\). On every deeper
layer where that tame inertia has been killed, the
construction gives a characteristic-zero integral
seed with residual datum \(S_F\). This is the
restriction of \(S\) with its transported frames.
The transported seed satisfies
\[
      S_F\in\mathcal A,\qquad
        -w_\beta,\,-w_\vartheta\in\operatorname{supp}(S_F).
\]
\end{established}
\begin{proof}
Proposition~\ref{prop:rational-rank-four-torus} gives the two
alternatives for the marked rational target in its normalized
basis. Lemma~\ref{lem:count-two-target-restriction} gives the restriction
calculation for these shapes: a nonzero normalized cocycle of
character \(\theta\) remains nonzero on the normal subgroup
containing a marked element \(\tau\) with \(\theta(\tau)\ne1\).
Indeed its vanishing there would give
\((1-\theta(\tau))f(g)=0\) for every \(g\) by conjugation.
Proposition~\ref{prop:count-three-fitted-layers} retains the
distinguishing elements here, and the same lemma treats the
transported local frame in the second shape.
For the other endpoint use Proposition~\ref{prop:count-three-tensor-endpoint}.
In the tensor basis its two marked factor entries supply
\(-w_\beta,-w_\vartheta\) and the prescribed nonzero local frame
entry. Killing the added tame inertia gives precisely the deep-layer
seed there, with the same marked reductions. These two calculations
apply to the two endpoints before any interpolating family is chosen.
\end{proof}

\proofstep{Simultaneous cup lifts.}
For support \(S_-\), row one of a product gives
\[
 (r(g)r(h))_{10}=r(g)_{10}d_0(h)+d_1(g)r(h)_{10}.
\]
Dividing by \(d_0(gh)\) yields the \(A\)-cocycle \(x\).
For \(S_+\), row two uses just columns zero and two,
and the same computation gives the \(B\)-cocycle \(y\). Normalization makes
these cocycles zero at the normalizer.
The first-column equations on the marked
decomposition groups give the local coboundaries
\[
 x(s)=(1-A(s))(N_v)_{10},\qquad
 y(s)=(1-B(s))(N_v)_{20}
\]
in the respective supports. In the first identity
row one of \(r(s)N_v\) uses only the first two
entries of its first column. In the second identity
row two uses only entries zero and two.

Apply the locally split cup lift of
\cite{count-two-cohomology}, with \(x\)
from \(S\) held fixed and a new \(B\)-cocycle \(y\)
chosen nonzero. Its finite coefficient module has
coordinates \((y,z')\), with action
\[
   (u,v)\longmapsto (B(g)u,(AB)(g)v+x(g)B(g)u).
\]
The result is a cochain \(z'\) with product
cross term \(x(g)B(g)y(h)\).
The cochain \(z=xy-z'\) has cross term
\(y(g)A(g)x(h)\), by the product rule for \(xy\).
Write \(\mathcal D=\operatorname{diag}(d_i)\) and
\[
 U(x,y,t)=
 \begin{pmatrix}
 1&0&0&0\\ x&1&0&0\\ y&0&1&0\\ t&y&-x&1
 \end{pmatrix}.
\]
At a marked place write \(x=(1-A)m_v\) and
\(y=(1-B)n_v\). For this local primitive the change \(z=xy-z'\) gives
\[
  z'=(1-AB)o'_v-xBn_v
     \quad\Longrightarrow\quad
  z=(1-AB)(m_vn_v-o'_v)-yAm_v .
\]
Here \(x=(1-A)m_v\) and \(y=(1-B)n_v\) are used on the same
decomposition group.

Consequently the matrices \(U(x,y,2z-xy)\mathcal D\)
with frames \(U(m_v,n_v,m_vn_v-2o'_v)\)
give a datum \(P_2\) of support \(S_0\).
In particular \(P_2\) has the
same \((1,0)\) coordinate as \(S\) in both global
matrices and local frames, and its root
\(-w_\Psi\) occurs. The new ramification is only
at finitely many extra split rational primes as
in the cup lift.

Next retain the \(B\)-cocycles of \(P_0\) and
\(P_2\) simultaneously. Use the direct sum of
the two cup coordinates with a common \(A\)-coordinate.
The finite action is
\[
 (u,v_0,v_2)\longmapsto
  (A(g)u,(AB)(g)v_0+y_0(g)A(g)u,
           (AB)(g)v_2+y_2(g)A(g)u).
\]
A cocycle \((x,v_0,v_2)\) for this action has the usual
\(A\)-cocycle law in its first coordinate and satisfies
\[
 v_i(gh)-v_i(g)-(AB)(g)v_i(h)=y_i(g)A(g)x(h),
                                  \qquad i=0,2 .
\]
These are the two cup equations with one common \(x\).
Prescribe the ramification direction \((1,0,0)\).
The duality argument gives a solution with \(x\ne0\)
and both lifts locally split at the old places. Hence it constructs two data \(P_3,P_4\)
with global and frame support in \(S_0\), using the
retained \((2,0)\) entries of \(P_0,P_2\),
respectively. They have the same \((1,0)\) entry in
each matrix and frame. Their common root
\(-w_\alpha\) occurs. The same local constants
are used for the common \(x\) at each marked place.

For every new prime, including \(w\), let \(C_w\) be
the degree-\(\ell\) cyclotomic subfield at that conductor
and set \(F_1=F_0\prod_w C_w\). The diagonal on each added
inertia group is trivial. For that inertia group \(I_w\), the cocycle factors as
\[
 I_w\longrightarrow I_w^{\rm tame,(\ell)}
 \simeq\Z_\ell(1)\longrightarrow\mathbf F_\ell(1)
 \longrightarrow \text{the residual coefficient module}.
\]
The degree-\(\ell\) cyclotomic field at \(w\) is totally ramified there, so its inertia subgroup lies in the kernel of the middle reduction.

On inertia at
\(w\) the tensor input will use that quotient as
well. Thus restriction removes the extra ramification
of all data at once. Over this tame \(F_1\), fix the
cohomological bound \(C\) for the growing cyclotomic
layers, and choose the layer \(F\) of the indicated
depth. The good target remains the restriction of
the same rational point. The identities of retained
entries and the nonvanishing needed for the links
persist on restriction with the transported flags.
Here the identities of global entries are
restrictions, and the nonzero root cocycles survive
by the \(\ell\)-power calculation of
\cite{count-two-cohomology}.

For the local frames one uses the identities
\[
 (N_v')_{i0} =
 \frac{(r(g)N_{v_0})_{i0}}{(r(g)N_{v_0})_{00}}
 \qquad (i=1,2)
\]
whenever the indicated conjugating matrix has unit
first pivot. This follows from the first column in
the lower--upper factorization defining the
transported frame. Each support is triangular in
a suitable coordinate order, so the principal pivots
are units and elimination keeps the same support. In both supports the denominator
equals \(d_0(g)\). For \(i=2\) with support \(S_+\),
row two uses only columns zero and two, and for
\(i=1\) with support \(S_-\), row one uses only
columns zero and one. This proves that the equalities
of retained frame coordinates survive transport.
The formula for \(i=1\) also transports the
nonzero seed frame coordinate to \(P_2\).
The rational lattice shape already supplies the
marked frame on the restricted target, so the
comparison uses its actual local flag.

\proofstep{The links and their endpoints.}
The path over \(F\) is
\[
                     P_0\longrightarrow P_3
                  \longrightarrow P_4
                  \longrightarrow P_2
                  \longrightarrow S_F .
\]
For the first and third links the two endpoints
have global and frame support in \(S_+\) and
the same \((2,0)\) data. Order the blocks as
\(\{0,2\},\{1,3\}\). The matrices are block lower
triangular in this order. Their first diagonal
block is fixed because it consists of the two
diagonal entries and the common \((2,0)\)
entry. In bases dual for the pairing, write \(A,D\) for the two diagonal blocks. The similitude equation in the paired blocks is
\[
 A^{\mathsf T}D=\mu I,\qquad D=\mu A^{-\mathsf T}.
\]
Thus fixing \(A\) and the similitude fixes \(D\). For the second
and fourth links the global support is \(S_-\)
and the frames are lower triangular.
Use blocks \(\{0,1\},\{2,3\}\), whose diagonal
blocks are fixed by the common \((1,0)\) data and
the similitude.

Thus at each link, for both the
representation and the frames, the difference
lies in the lower-left block ideal of a block
triangular algebra. This ideal is square-zero
and is preserved by the symplectic adjoint in
each order because the two blocks are paired.
In that algebra let \(A_s=(1-s)A_0+sA_1\) and
\(B_s=(1-s)B_0+sB_1\) be matrices built from two endpoint
pairs. Their differences have zero product, giving
\[
 A_s B_s
      =(1-s)A_0B_0+sA_1B_1,
 \qquad
 (A_s)^\dagger=(1-s)A_0^\dagger+sA_1^\dagger .
\]
For each link write its global matrices \(r_s\) and local
frames \(N_{v,s}\) by affine interpolation. The two product
identities yield
\[
 r_s(g)r_s(h)=r_s(gh),\qquad
 r_s(g)^\dagger r_s(g)=\nu(g)I,\qquad
 r_s(h)N_{v,s}=N_{v,s}B_{v,s}.
\]
In the last identity \(h\) is local and \(B_{v,s}\) is the
interpolation of the conjugated upper triangular matrices.
It retains the same ordered diagonal. In particular the normalization, prescribed
unramified quotients and residual diagonal hold
on all four polynomial lines.

The first bridge point \(P_3\) has the same
\((2,0)\) data as \(P_0\) and a nonzero
\(-w_\alpha\) root. Thus it has the first type
of ordinary good residual datum if the target
has global \(-w_\Psi\), and the frame type if
the target has only the local \((2,0)\) data.
For \(P_4\) the two nonzero roots are explicit:
\[
 \{-w_\alpha,-w_\Psi\}\subseteq\operatorname{supp}_{\rm glob}(P_4).
\]
This verifies the first ordinary test using the global matrices alone.

Write \(r_2\) and \(N_{2,v}\) for the global matrices and marked frames of \(P_2\).
The \((1,0)\) entry agrees with the seed. In support \(S_0\) the symplectic adjoint pairs it with \((3,2)\), while the local frames retain their nonzero seed entry. Thus
\[
 (r_2(g))_{10}=(r_2(g))_{32}=0\quad\text{for every }g,
 \qquad (N_{2,v})_{10}\ne0\quad\text{at every marked }v.
\]
It follows that \(P_2\in\mathcal A\); its root \(-w_\Psi\) occurs. The endpoint
\(S_F\) is in \(\mathcal A\) with both
\(-w_\beta,-w_\vartheta\) occurring. The last
two good-point conditions are the extra cases
required for the rank-four case.
\label{cor:ordinary-count-iii-chain-1}

\begin{construction}
\(\displaystyle
P_0 \xrightarrow[\ (2,0)\ ]{S_+} P_3
    \xrightarrow[\ (1,0)\ ]{S_-} P_4
    \xrightarrow[\ (2,0)\ ]{S_+} P_2
    \xrightarrow[\ (1,0)\ ]{S_-} S_F
\)
\caption{The four affine residual lines. Above each arrow is the allowed global support on that line, below is the entry shared by the two endpoints in both the global and frame data. Fixing the paired diagonal blocks makes the interpolation multiplicative.}\label{fig:four-residual-line-supports}
\end{construction}

\proofstep{Descent of the residual path.}

The links have been verified after making coefficient choices for their cocycles. They enter the component argument only once all four descend to the same finite residue field.

The matrices in a polynomial line are locally
constant over the discrete ring \(\overline{\mathbf F}_\ell[X]\).
Each continuous polynomial-valued matrix has finite image by
compactness and discreteness. Let \(S\) be the union of all coefficients of those images and frames across the four lines. Both the set of matrices and their polynomial degrees are finite, so
\[
 k'=\mathbf F_\ell(S),\qquad [k':\mathbf F_\ell]<\infty .
\]
All four matrix families have entries in \(k'[X]\).
The algebraic equations, inertia equations and
nonvanishing conditions can be read there
because scalar extension is injective. The four
lines and their good endpoints consequently
descend to a single finite residue field. This
extension can be made together with the
coefficient extensions for the actual
characteristic-zero seed, the small target model
and the two ordinary levels.
\label{cor:ordinary-count-iii-chain-2}
\end{proof}
